% !TeX root = ../MAgent_Theory_2.0_companion.tex
\chapter{Links and Physical Geometry}
\label{comp:geometry}

The agent graph and the context base describe different directions of comparison. An edge $j\to i$ compares agents at one context. A path in the base $\mathcal C$ compares states at successive contexts. A product of communication links can be nontrivial without defining curvature on $\mathcal C$. Conversely, one agent can carry a curved connection over contexts; collective averaging is not required.

\section{Relating neighboring contexts}

To construct that connection, suppress the sector label and fix one agent. On a coordinate neighborhood, take a twice continuously differentiable presentation $U(c)$ and a continuously differentiable matrix-valued one-form $a=a_\mu(c)dc^\mu$. Declare directional links whose first-order behavior is
\begin{equation}
\Theta(c+\delta c,c)
=U(c+\delta c)e^{-a_\mu(c)\delta c^\mu}U(c)^{-1}
+O(|\delta c|^2).
\end{equation}
The displacement enters linearly at first order, and zero displacement gives identity. These are additional properties beyond having a finite interaction matrix at each context. For a principal $G$-connection, the generators must belong to the represented Lie algebra; unrestricted matrices instead define general linear bundle connections.

Using $D=d+B$, expansion identifies
\begin{equation}
B=UaU^{-1}-dUU^{-1}.
\end{equation}
The first term carries the declared directional interaction into displayed coordinates. The second compensates for the change of presentation along the path. Their relationship is exactly the operator identity
\begin{equation}
d+B=U(d+a)U^{-1}.
\end{equation}
Indeed, acting on a vector field $f$ gives $U d(U^{-1}f)=df-dUU^{-1}f$; the remaining term is $UaU^{-1}f$. Squaring the operator therefore gives
\begin{equation}
F=dB+B\wedge B=U(da+a\wedge a)U^{-1}.
\end{equation}
This is the mathematical connection to gauge theory, with the sign convention fixed by the displayed covariant derivative \citep{Tong2018}.

\section{Curvature and paths}

Curvature need not come from spatially varying coefficients. On a two-dimensional context patch choose $U=I$, $a_x=E_{12}$, and $a_y=E_{21}$. Both directional generators have rank one, but
\begin{equation}
F_{xy}=[E_{12},E_{21}]
=\begin{pmatrix}1&0\\0&-1\end{pmatrix}.
\end{equation}
The matrix $E_{ab}$ has one unit entry in position $(a,b)$. Applying the two infinitesimal operations in opposite orders gives different results. This example supplies directional data; it does not derive them from the agents' laws. A presentation change alone, $B=-dUU^{-1}$, instead has zero curvature.

Finite path transport solves
\begin{equation}
\frac{dT}{dt}=-B(\dot\gamma(t))T(t),\qquad T(0)=I.
\end{equation}
Its solution is a path-ordered exponential. Concatenating paths composes their transports; reversing the same path gives the inverse. Two independently specified attending edges can have unrelated generators, so their opposite labels do not impose this inverse relation. The distinction is necessary when borrowing Wilson-loop language from gauge theory \citep{Wilson1974}.

\section{Local descriptions and global compatibility}

Endpoint covariance requires
\begin{equation}
\Theta'(c_2,c_1)=g(c_2)\Theta(c_2,c_1)g(c_1)^{-1}.
\end{equation}
Differentiating the endpoints produces the affine law
\begin{equation}
B'=gBg^{-1}-dg g^{-1}.
\end{equation}
Likewise, keeping the intrinsic transport fixed under $U'=Uh$ requires $a'=h^{-1}ah+h^{-1}dh$. A directional receiving--sending product that transforms only by conjugation is a tensorial one-form. It can supply the difference between two connections, but does not by itself supply the full affine transformation law.

Global compatibility adds another requirement. If patch coordinates obey $v_\alpha=g_{\alpha\beta}v_\beta$, then $g_{\alpha\beta}g_{\beta\gamma}=g_{\alpha\gamma}$ on triple overlaps, and local connection forms must satisfy the affine law there. A nontrivial bundle may have no single global presentation. Even a trivial bundle can carry flat transport with global memory: on a circle, the scalar connection $B=b d\theta$ has $F=0$ but holonomy $e^{-2\pi b}$ around one turn. For real $b\ne0$, this differs from identity.

Physical interpretation requires a model of the contexts, an agent mechanism producing compatible directional links, and observable responses to their transport. A connection alone does not supply spacetime or quantum dynamics.
